
Hoved og under spørgsmål:

1. Hvordan kan SciML anvendes  til modellering af HESEL ligningsystemet ?
Her vil vi bruge PINNs til at modellere HESEL-ligningsystemet og sammenligne det med numeriske simuleringer.

2. Hvad er effektten af at træne på ligningerne hver for sig i modsetning til det fulde ligningsystem?
Hvilken effekt har det at træne på lav resulition data i forhold til høj resulition data? 

3. Kan det generalisere til usete træmings prunkter og flere datasæt/begyndelsesbetingelser?
Design of experiments, bayesisk optimering

4. Kan vi effektivsiere modellen ?
via pruning, lav rank/ reduced order modelling, eller andre metoder?
